%% kapa: exam=1 round=27 year=2016 period=2 form=A subject=appraisal_law subject_ko=감정평가관계법규 seq=25 number=25
%% source: Q-Net 감정평가사 기출문제 (cst00309) artlSeq=5169402 file=extracted/1008/5169402/제27회 감정평가사 1차 문제/제27회 감정평가사 2교시문제(A형).pdf
%% conversion: kordoc 4.24.51 -> tools/gen1.py (auto); review=proofread-sonnet; flags=none
\begin{question}{25}
건축법령상 용어에 관한 설명으로 옳지 않은 것은?
\begin{choices}
  \item ‘지하층’이란 건축물의 바닥이 지표면 아래에 있는 층으로서 바닥에서 지표면까지 평균높이가 해당 층 높이의 2분의 1 이상인 것을 말한다.
  \item 건축물을 이전하는 것은 ‘건축’에 해당하지 않는다.
  \item ‘리모델링’이란 건축물의 노후화를 억제하거나 기능 향상 등을 위하여 대수선하거나 일부 증축하는 행위를 말한다.
  \item 층수가 25층이며, 높이가 120미터인 건축물은 ‘고층건축물’에 해당한다.
  \item 피뢰침은 ‘건축설비’에 해당한다.
\end{choices}
\end{question}
\answer{2}
